\chapter{Modalidades de TCC}

Os cursos podem admitir diferentes modalidades de Trabalho de Conclusão de Curso, conforme o projeto pedagógico e o regulamento específico. A presença de uma modalidade neste manual não constitui autorização automática: a decisão final cabe ao curso.

\regracurso

\section{Modalidades admitidas}

As modalidades a seguir são referências para planejamento. Em todas elas, o trabalho deve demonstrar problema ou demanda claramente definidos, fundamentação, método, análise e contribuição acadêmica compatíveis com a formação do curso.

\subsection{Monografia}

A monografia é adequada quando a investigação exige desenvolvimento extenso da fundamentação, do método e da análise. Pode ser bibliográfica, documental, empírica, aplicada ou mista.

Sua estrutura costuma reunir elementos pré-textuais, introdução, referencial teórico ou revisão de literatura, metodologia, resultados e discussão, parte conclusiva, referências e, quando necessários, apêndices e anexos. Os títulos e a divisão do desenvolvimento devem acompanhar o objeto e a tradição da área.

\subsection{Artigo científico}

O artigo científico exige recorte preciso, unidade e redação concisa. Pode ser adotado quando o curso pretende aproximar o estudante da comunicação em periódicos ou eventos, sem reduzir o rigor esperado do TCC.

Sua organização costuma incluir título, autoria, resumo, palavras-chave, introdução, método, resultados, discussão, conclusão e referências. Resultados e discussão podem aparecer juntos ou separados, conforme a área e as regras do curso ou do veículo escolhido.

Quando o artigo constituir o formato de entrega, devem ser observadas a ABNT NBR 6022:2018 e as regras do curso. A eventual adoção das instruções editoriais de um periódico não dispensa os elementos institucionais exigidos para depósito, salvo autorização expressa.

\subsection{Relatório técnico ou científico}

O relatório é adequado para intervenção, diagnóstico, prática profissional, avaliação institucional, desenvolvimento tecnológico, processo ou serviço. Ele não deve se limitar a narrar atividades.

O texto deve apresentar contexto, problema, objetivos, procedimentos, resultados, análise, recomendações quando cabíveis, parte conclusiva e referências. A descrição do que foi feito precisa ser acompanhada da justificativa das escolhas e da avaliação do que foi alcançado.

Quando aplicável, a apresentação do relatório deve considerar a ABNT NBR 10719:2015, além das exigências do PPC ou regulamento.

\subsection{Produto técnico, tecnológico ou educacional}

Essa modalidade pode envolver protocolo, aplicativo, material didático, guia, cartilha, instrumento, oficina, proposta de intervenção ou outra solução aplicada. O produto deve ser acompanhado de texto acadêmico que explique o problema, o público, os fundamentos, os requisitos, o processo de desenvolvimento, a validação e os limites de aplicação.

\subsection{Portfólio acadêmico ou de projetos}

Quando previsto pelo curso, o portfólio pode reunir produções articuladas que evidenciem percurso formativo, decisões, resultados e reflexão crítica. Ele não é uma simples coleção de arquivos: deve possuir critérios de seleção, organização, contextualização e análise das evidências apresentadas.

\subsection{Revisão de literatura}

A revisão pode constituir modalidade de TCC quando o objetivo é mapear, sintetizar ou analisar criticamente a produção existente. Deve explicitar bases consultadas, estratégia de busca, período, critérios de inclusão e exclusão e procedimento de análise. Uma sequência de resumos sem método e sem síntese crítica não caracteriza revisão acadêmica suficiente.

\subsection{Estudo de caso}

O estudo de caso pode ser admitido como formato de trabalho ou empregado como estratégia de pesquisa em outra modalidade. Deve delimitar o caso, justificar sua escolha, identificar as fontes de evidência e desenvolver análise que ultrapasse o relato descritivo.

\section{Escolha da modalidade}

A modalidade deve ser escolhida com o orientador e confirmada conforme o regulamento do curso. A decisão deve considerar a pergunta de pesquisa, os objetivos, a natureza dos resultados, o tempo disponível, o acesso a dados, as exigências éticas e a forma de avaliação. A modalidade organiza a entrega acadêmica; ela não substitui o método da pesquisa.

\caixaatencao{Pôster e slides são formas de apresentação do trabalho, não modalidades do documento de TCC.}

As orientações para pôsteres e slides devem ser publicadas como material de apoio à apresentação. Caso a UnDF adote pôster científico institucional, a norma específica vigente deve ser conferida antes da publicação do modelo.
